% !TeX program = pdflatex
\documentclass[10pt,a4paper]{article}

\usepackage[utf8]{inputenc}
\usepackage[T1]{fontenc}
\usepackage{geometry}
\usepackage[hidelinks]{hyperref}
\usepackage{xcolor}

\geometry{
  top=1.7cm,
  bottom=1.8cm,
  left=1.8cm,
  right=1.8cm
}

\pagestyle{empty}
\frenchspacing
\setlength{\parindent}{0pt}
\setlength{\parskip}{0.45em}

\definecolor{mydarkgreen}{RGB}{0,0,0}

\begin{document}

\begin{minipage}[t]{0.48\textwidth}
\textbf{\large Javier Andres TARAZONA JIMENEZ}\\
Étudiant ingénieur IA \\
Machine Learning, Vision 3D, LLMs et prototypage \\
Massy, Île-de-France, France\\
\href{mailto:javitar06@gmail.com}{javitar06@gmail.com}\\
+33 07 46 36 97 75\\
\href{https://www.linkedin.com/in/javier-andres-tarazona-jimenez-84b489222/}{LinkedIn}
\end{minipage}
\hfill
\begin{minipage}[t]{0.42\textwidth}
\raggedleft
Dassault Systèmes \\
Service Recrutement \\
Équipe Generative Experiences -ENOVIA  \\
Campus de Vélizy-Villacoublay \\
Vélizy-Villacoublay, France \\
\end{minipage}

\vspace{1.1cm}

\raggedleft
Palaiseau, le 24 juin 2026

\vspace{0.7cm}

\raggedright
\textbf{Objet : Candidature à l’apprentissage Ingénieur Intelligence Artificielle – ENOVIA / Generative Experiences – Réf. 547433}

\vspace{0.5cm}

{\small

Madame, Monsieur,

Actuellement étudiant colombien en double diplôme entre l’Universidad Nacional de Colombia et ENSTA Paris, au sein de l’Institut Polytechnique de Paris, je souhaite vous présenter ma candidature pour l’apprentissage Ingénieur Intelligence Artificielle (F/H), au sein de l’équipe « Generative Experiences » d’ENOVIA chez Dassault Systèmes, à partir de septembre 2026.

J’ai choisi l’informatique par intérêt pour la compréhension des systèmes complexes, par goût pour la construction de solutions concrètes à partir d’idées abstraites, et par volonté d’acquérir une autonomie technique solide. L’intelligence artificielle s’est ensuite imposée comme une évolution naturelle de ce parcours : passer de la programmation de systèmes à la conception de systèmes capables de percevoir, de raisonner, d’apprendre et d’assister la décision. Ce domaine m’attire parce qu’il combine rigueur mathématique, expérimentation et problèmes ouverts. La vision par ordinateur et les agents fondés sur de grands modèles de langage m’intéressent particulièrement, car ils relient perception, données, raisonnement et interaction utilisateur dans des systèmes utiles.

L’intérêt que je porte à Dassault Systèmes dépasse largement l’image d’une entreprise de logiciels 3D. Je vois surtout en vous un acteur majeur des jumeaux virtuels, de la simulation, de la géométrie numérique et des plateformes industrielles collaboratives. La plateforme 3DEXPERIENCE et les applications ENOVIA correspondent à une vision de l’IA que je trouve particulièrement stimulante : connecter les personnes, les idées, les données et les processus afin d’améliorer la conception, la collaboration et la prise de décision. Pour un profil comme le mien, intéressé par une IA ayant un impact économique et technologique réel, ce cadre est plus motivant qu’une application isolée. Il correspond aussi à ma manière de réfléchir : comprendre les interactions entre les composants d’un système, modéliser leurs contraintes et construire des solutions utiles dans un environnement complexe.

Votre offre présente un noyau qui correspond directement à mon parcours : conception de solutions d’IA prédictive ou générative, développement d’assistants virtuels multi-agents avec LangChain et LangGraph, stratégies de résolution de problèmes comme ReAct ou Plan&Execute, collecte de données, évaluation, documentation et démonstration de la valeur produite. Mon stage actuel à ENSTA Paris – U2IS me permet de travailler sur la vision 3D auto-supervisée, la synthèse de nouvelles vues, la génération de données 3D synthétiques avec Blender, les rendus multi-vues, les métadonnées de caméras et le pré-entraînement avec PyTorch et Vision Transformers. Cette expérience m’a appris à construire des pipelines expérimentaux, préparer des données, comparer des configurations et documenter des résultats de manière rigoureuse.

Je peux également apporter une base solide en développement logiciel, en vision par ordinateur et en prototypage d’applications d’IA. Chez APEDYS91, j’ai travaillé sur un pipeline NLP combinant règles, LLMs génératifs, contrôle d’entités et sélection dynamique de modèles. Chez Engin A.I, j’ai contribué à l’analyse d’images, à l’optimisation de pipelines OpenCV/NumPy et à la structuration de bibliothèques Python en architecture modulaire. Dans mes projets académiques, j’ai développé des systèmes de détection, de segmentation, de suivi visuel et des API REST avec FastAPI. Cette double compétence « IA appliquée et ingénierie logicielle » me semble essentielle pour transformer une idée de recherche en prototype fonctionnel, testable et progressivement industrialisable.

Rejoindre Dassault Systèmes représenterait pour moi une transition structurante entre l’IA académique et l’IA industrielle. J’y vois l’opportunité d’apprendre comment des modèles, des agents, des algorithmes et des prototypes deviennent des composants intelligents intégrés à des plateformes utilisées par de grandes industries. Ce cadre me permettrait de faire évoluer mon profil vers une expertise plus complète : IA générative, agents LLM, vision 3D, segmentation, géométrie, ingénierie système et produit réel.

Je serais heureux de pouvoir échanger avec vous afin de vous présenter plus en détail ma motivation et l’adéquation de mon parcours avec vos besoins.

\textbf{Cordialement},

Javier Andres TARAZONA JIMENEZ


}

\end{document}